\section{Introduction}

% prediction-based compression重要
The volume of scientific data generated by high-performance computing (HPC) simulations and scientific instruments continues to grow rapidly, placing increasing pressure on storage systems, I/O bandwidth, and data movement~\cite{lockwood2017storage,huang2023towards,lewis2025machine,giannakou2024understanding,lockwood2019understanding}. Error-bounded lossy compression has emerged as an important technique for alleviating these challenges~\cite{cappello2019use,di2025survey,cappello2025support}, as it substantially reduces data volume while preserving user-specified accuracy~\cite{lindstrom2014fixed,di2016fast,huang2023cuszp,gong2023mgard,li2023lossy}.

% 现有方法中的tuning
Achieving strong compression performance requires more than finding better compression algorithm. It also depends critically on how the compressor is configured for the target data and the user’s compression objective. These configurations may include not only parameter values, but also which compression components are executed and how they are organized. Consequently, different configurations of the same compressor can exhibit substantially different compression ratios and data distortion. For example, QoZ~\cite{liu2022dynamic,liu2024high}, a prediction-based compressor, dynamically selects among candidate predictors during online compression according to the input data. At the compressor-design level, transform-based compressors such as ZFP~\cite{lindstrom2014fixed,diffenderfer2019error} have continued to evolve through refinements to their transform designs and coding strategies. At an even coarser level, some compression frameworks~\cite{liang2019significantly} combine multiple existing compressors and select which compressor to invoke for a particular dataset or data field.

% 现有方法中tuning的问题
Despite the strong performance achieved by existing compressors and tuning strategies, further improvements may be possible through a larger and more systematic exploration of the configuration space. On the one hand, candidate configurations are commonly evaluated by executing the full compression pipeline on sampled data~\cite{liang2022sz3,zhao2021optimizing}, which provides reliable rate–distortion measurements but limits the number of candidates that can be explored. Prior studies have attempted to reduce this cost by skipping parts of the compression pipeline and using estimated metrics or lightweight proxies to guide configuration selection~\cite{liang2019significantly,jin2022improving,liu2022dynamic,liu2021improving,liu2024high}. Although these approaches can substantially reduce tuning overhead, the resulting estimates are not always sufficiently reliable to distinguish among candidate configurations. On the other hand, some configuration dimensions that may offer additional compression benefits have received limited attention or have not yet been systematically explored.

% tuning框架
To address these challenges, we propose a hierarchical tuning framework that determines when and how each configuration dimension should be explored based on its potential performance benefit, the scope over which the tuning decision can be reused, and the cost and fidelity of available evaluation methods. Configuration dimensions whose optimal choices generalize well across datasets, or whose expected performance gains are insufficient to justify repeated tuning costs, can be studied offline and fixed by compressor developers. Configurations that depend on dataset-level characteristics can be tuned by users on representative data and reused across related files, while highly data-dependent configurations are tuned online during compression using low-cost evaluation methods. By assigning different configuration dimensions to appropriate tuning stages, the framework enables broader configuration exploration without repeatedly incurring prohibitive tuning overhead.

% asz
To demonstrate this framework, we develop ASZ as an instantiation for prediction-based error-bounded compression. At design time, ASZ studies broadly applicable algorithmic structures and control parameters whose benefits do not justify repeated runtime tuning. At the dataset level, it adaptively selects an ordered sequence of prediction dependency patterns and tunes the data traversal order used in each prediction round. At the block level, ASZ performs large-scale online tuning of predictor coefficients. To make this fine-grained exploration practical, we develop a tailored entropy–distortion evaluation criterion and an optimized evaluation procedure that reuses data-dependent statistics across coefficient candidates, thereby avoiding repeated prediction and full lossless encoding.

% 测试结果
Across eight evaluated scientific datasets, ASZ consistently achieves better rate–distortion performance than QoZ~\cite{liu2022dynamic,liu2024high}, a state-of-the-art prediction-based compressor, while maintaining comparable running time. Although ASZ incurs an additional preprocessing cost, the resulting configurations can be reused across related files or timesteps, allowing the overhead to be amortized in practical scientific data management workflows. ASZ also outperforms leading non-prediction-based compressors on most of the evaluated datasets.

Our main contributions are summarized as follows:

\begin{itemize}
\item We propose a hierarchical tuning framework that systematically organizes configuration exploration across offline, dataset-level, and online stages.
\item We identify and explore previously underexamined configuration dimensions in prediction-based compression, including flexible ordering of prediction dependency patterns and tunable data traversal orders. We further enable efficient block-level coefficient tuning through a geometric $H_1$--PSNR criterion and reusing coefficient-independent terms in residual evaluation.
\item Based on this framework, we develop ASZ. Experiments on eight scientific datasets show that ASZ achieves better rate--distortion performance than QoZ with comparable runtime, while outperforming leading non-prediction-based compressors on most datasets.
\end{itemize}

The rest of this paper is organized as follows. Section~\ref{sec:related work} reviews related work on prediction-based error-bounded lossy compression and predictor selection methodologies used in existing lossy compression frameworks. Section~\ref{sec:preliminary study} includes a preliminary study of stencil patterns and provides some important takeaways. Section~\ref{sec:asz} presents a detailed analysis of entropy-distortion-based bitrate estimation, describes our fast predictor evaluation method, and introduces the resulting compression algorithm, ASZ. Section~\ref{sec:performance} evaluates ASZ against state-of-the-art prediction-based compressors on real-world datasets from different domains. Section~\ref{sec:conclusion} concludes the paper.